\uuid{2Fp1}
\titre{Points stationnaires}
\auteur{hanani ; mzari}
\organisation{exo7}
\licence{CC BY-NC-SA 4.0}
\datecreate{2020-01-20}
\niveau{L1}
\module{Géométrie}
\chapitre{Courbes planes}
\sousChapitre{Courbes paramétrées}
\difficulte{3}
\theme{Courbes paramétrées, points stationnaires}
\exo7qcmid{lille-2.courbes.006}

\contenu{
\begin{question}
\qclassification{244.01}
\qtags{motcle=points stationnaires, branches infinies}
La trajectoire $\Gamma$ d'une particule en mouvement est donnée par les équations
$$x=\frac{1}{2}(t^2-2t)\quad \mbox{et} \quad  y=\frac{1}{3}t^{3}-\frac{1}{2}t^2.$$
\begin{answers}
\bad{Le point de paramètre $1$ est un point de rebroussement de seconde espèce.}
\good{La tangente à $\Gamma$ au point $(x(1),y(1))$ est la droite d'équation $y=x$.}
\bad{Le point de paramètre $0$ est un point stationnaire.}
\good{$\Gamma$ admet une branche parabolique de direction asymptotique l'axe des $y$.}
\end{answers}
\begin{explanations}
Avec $f(t)=(x(t),(y(t))$. On a : $f'(1)=(0,0)$, $f''(1)=(1,1)$ et $f^{(3)}(1)=(0,2)$. Le point de paramètre $1$ est un point de rebroussement de première espèce et la tangente en ce point est la droite d'équation $x-y=0$. Enfin
$$\displaystyle \lim _{t\to \pm \infty}x(t)
=+\infty,\qquad \lim _{t\to \pm \infty}y(t)=\pm\infty\quad\mbox{et}\quad \lim _{t\to \pm \infty}\frac{y(t)}{x(t)}=\pm\infty .$$
Donc $\Gamma$ admet une branche parabolique de direction asymptotique l'axe des ordonnées.
\end{explanations}
\end{question}
}
